\documentclass[10pt,letterpaper]{article}

\usepackage[letterpaper,margin=0.74in,top=0.62in,bottom=0.68in]{geometry}
\usepackage[T1]{fontenc}
\usepackage[utf8]{inputenc}
\usepackage[sfdefault]{sourcesanspro}
\renewcommand{\familydefault}{\sfdefault}
\usepackage{microtype}
\usepackage{array}
\usepackage{tabularx}
\usepackage{enumitem}
\usepackage{titlesec}
\usepackage{needspace}
\usepackage{xcolor}
\usepackage{xurl}
\usepackage{hyperref}
\usepackage{fancyhdr}
\usepackage{lastpage}

\input{glyphtounicode}
\pdfgentounicode=1

\definecolor{accent}{HTML}{163A5F}
\definecolor{muted}{HTML}{555555}
\definecolor{rulegray}{HTML}{B7C0C9}

\hypersetup{
  colorlinks=true,
  urlcolor=accent,
  linkcolor=accent,
  pdftitle={Curriculum Vitae - Yifan Xiong},
  pdfauthor={Yifan Xiong},
  pdfsubject={Academic curriculum vitae}
}
\urlstyle{same}

\pagestyle{fancy}
\fancyhf{}
\renewcommand{\headrulewidth}{0pt}
\renewcommand{\footrulewidth}{0pt}
\fancyfoot[C]{\color{muted}\scriptsize Yifan Xiong \enspace | \enspace Curriculum Vitae \enspace | \enspace Page \thepage\ of \pageref*{LastPage}}

\setlength{\parindent}{0pt}
\setlength{\parskip}{0pt}
\setlength{\tabcolsep}{0pt}
\setlength{\emergencystretch}{2em}
\raggedright

\titleformat{\section}
  {\large\bfseries\color{accent}}
  {}{0pt}{}
  [\vspace{0.12em}{\color{rulegray}\titlerule[0.4pt]}]
\titlespacing*{\section}{0pt}{0.95em}{0.52em}

\newcommand{\cventry}[4]{%
  \needspace{3.2\baselineskip}%
  \begin{tabularx}{\textwidth}{@{}X@{\hspace{0.35em}}>{\raggedleft\arraybackslash}p{1.70in}@{}}
    \textbf{#1} & \textbf{#2} \\
    \textit{#3} & \textit{#4}
  \end{tabularx}%
}

\newcommand{\projectentry}[3]{%
  \needspace{5.2\baselineskip}%
  \begin{tabularx}{\textwidth}{@{}X@{\hspace{0.35em}}>{\raggedleft\arraybackslash}p{1.70in}@{}}
    \textbf{#1} & \textbf{#2} \\
    \multicolumn{2}{@{}l@{}}{\textit{#3}}
  \end{tabularx}%
}

\newenvironment{cvitems}{%
  \begin{itemize}[
    leftmargin=1.2em,
    label=\raisebox{0.18ex}{\tiny$\bullet$},
    itemsep=2.2pt,
    parsep=0pt,
    topsep=3pt,
    partopsep=0pt
  ]
}{\end{itemize}\vspace{0.3em}}

\newcommand{\skillrow}[2]{\textbf{#1:} #2\par}
% Use a raised math-font dagger for a clear, sword-like equal-contribution mark.
\newcommand{\equalcontrib}{\raisebox{0.6ex}{\scriptsize\ensuremath{\dagger}}}

\begin{document}

\fontsize{10.5}{12.4}\selectfont

\begin{center}
  {\fontsize{23}{25}\selectfont\bfseries Yifan Xiong}\par
  \vspace{4pt}
  \small
  Urbana-Champaign, Illinois \enspace | \enspace
  \href{mailto:yifanx14@illinois.edu}{yifanx14@illinois.edu} \enspace | \enspace
  \href{https://yifanxiong272.github.io/}{yifanxiong272.github.io}
  \\
  \href{https://github.com/yifanxiong272}{GitHub} \enspace | \enspace
  \href{https://scholar.google.com/citations?hl=en\&user=xyP48YoAAAAJ}{Google Scholar}
\end{center}

\vspace{-0.2em}

\section{Education}

\cventry
  {University of Illinois Urbana-Champaign}
  {Sep. 2026 - Present}
  {Ph.D. Student in Computer Science; Advisor: \href{https://yilinglou.github.io/}{Prof. Yiling Lou}}
  {Urbana-Champaign, IL}

\vspace{0.35em}

\cventry
  {Nanjing University}
  {Sep. 2021 - Jul. 2025}
  {B.Eng. in Software Engineering; GPA: 4.54/5.00}
  {Nanjing, China}

\section{Research Interests}

Software engineering, formal methods, and programming languages, with a focus on software testing, formal and runtime verification, program analysis, and their intersections with AI and agentic systems.

\section{Publications}

\begin{enumerate}[leftmargin=1.45em,label={\arabic*.},itemsep=6pt,topsep=1pt,parsep=0pt]
  \item Jinglin Dai\equalcontrib, \textbf{Yifan Xiong}\equalcontrib, Lezhi Ma, Shangqing Liu, and Lei Bu. ``Temporal Specification Oriented Fuzzing for Trigger-Action-Programming Smart Home Integrations.'' \textit{Proceedings of the 48th International Conference on Software Engineering (ICSE 2026), Research Track}. \href{https://doi.org/10.1145/3744916.3764543}{DOI} | \href{https://sites.google.com/view/hafuzz}{Project}

  \item Ruixin Qin\equalcontrib, \textbf{Yifan Xiong}\equalcontrib, Yifei Lu, and Minxue Pan. ``Demystifying and Assessing Code Understandability in Java Decompilation.'' \textit{arXiv preprint arXiv:2409.20343}, 2024. \href{https://arxiv.org/abs/2409.20343}{arXiv}

  \item Chun Li, \textbf{Yifan Xiong}, Zhong Li, Wenhua Yang, and Minxue Pan. ``Mobile Test Script Generation from Natural Language Descriptions.'' \textit{2023 IEEE 23rd International Conference on Software Quality, Reliability, and Security (QRS)}, pp. 348--359, 2023. \href{https://doi.org/10.1109/QRS60937.2023.00042}{DOI}
\end{enumerate}

{\footnotesize\equalcontrib\ Equal contribution.}

\section{Academic Positions}

\cventry
  {University of Illinois Urbana-Champaign}
  {Fall 2026}
  {Teaching Assistant, \href{https://siebelschool.illinois.edu/academics/courses/cs527}{CS 527: Software Testing and Analysis}}
  {Urbana-Champaign, IL}

\vspace{0.35em}

\cventry
  {Software Engineering Group, Nanjing University}
  {Aug. 2025 - Dec. 2025}
  {Research Assistant}
  {Nanjing, China}

\section{Research Experience}

\projectentry
  {HAFuzz: Temporal Specification Oriented Fuzzing for Smart Home Integrations}
  {Nov. 2023 - Oct. 2025}
  {Research Assistant; Advisor: Prof. Lei Bu, Nanjing University}
\begin{cvitems}
  \item Led the full implementation, including modeling IFTTT-style IoT systems, developing a specialized runtime-verification fuzzer, designing and refining seed-selection algorithms, and improving overall efficiency.
  \item Organized and conducted experiments comparing our tool with model-checking-based approaches on a comprehensive, carefully curated dataset, demonstrating improvements in efficiency and accuracy.
\end{cvitems}

\projectentry
  {Extending NNSmith with LLM-Generated API Constraints}
  {Jul. 2024 - Nov. 2024}
  {Research Assistant; Advisor: Prof. Lingming Zhang, University of Illinois Urbana-Champaign}
\begin{cvitems}
  \item Extended NNSmith's API-constraint workflow with LLM-generated candidates, expanding the set of supported constraints and the diversity of generated deep-learning models.
  \item Developed augmentation strategies to improve the accuracy of generated API constraints.
  \item Designed validation, repair, and coverage-measurement components; used the extended workflow to identify real-world bugs in deep-learning libraries.
\end{cvitems}

\newpage

\projectentry
  {Assessing Code Understandability in Java Decompilation}
  {May 2023 - Jun. 2024}
  {Research Assistant; Advisor: Prof. Minxue Pan, Nanjing University}
\begin{cvitems}
  \item Built a SonarQube-based analysis workflow and evaluated source and decompiled Java code across diverse projects and mainstream decompilers, including Jadx, FernFlower, and CFR.
  \item Trained an n-gram language model, conducted manual assessments using established Java conventions, and identified recurring decompilation patterns that reduce readability.
  \item Helped design Cognitive Complexity$^{D}$, a decompilation-aware metric that more accurately captures understandability differences between source and decompiled code than conventional metrics.
\end{cvitems}

\projectentry
  {Reinforced-Attention Policy for LibFewShot}
  {Sep. 2023 - Jan. 2024}
  {Research Project}
\begin{cvitems}
  \item Implemented a reinforced-attention framework combining few-shot learning with Markov decision processes and adapted the model to limited hardware resources.
  \item Evaluated multiple backbone networks and meta-learners in LibFewShot, successfully reproducing the reported improvements across varied configurations.
\end{cvitems}

\projectentry
  {GenDroid: Mobile Test Script Generation from Natural Language}
  {Sep. 2022 - Mar. 2023}
  {Research Assistant; Advisor: Prof. Minxue Pan, Nanjing University}
\begin{cvitems}
  \item Combined SentenceBERT features with a random forest classifier to match natural-language test intents to Android UI widgets; contributed path planning over DroidBot-generated UI transfer graphs.
  \item Identified omitted actions and application-state transitions to make multi-step test execution more robust and reduce cascading failures.
  \item Helped design and run experiments demonstrating improved intent coverage and widget-matching accuracy across diverse Android applications.
\end{cvitems}

\section{Honors and Awards}

\begin{tabularx}{\textwidth}{@{}X >{\raggedleft\arraybackslash}p{0.9in}@{}}
  Outstanding Paper Award, 18th National College Students Innovation Annual Conference, China & 2025 \\
  First-Class People's Scholarship, Nanjing University & 2023 - 2025
\end{tabularx}

\section{Technical Skills}

\skillrow{Programming}{Python, Java, C/C++, Bash, SMV, LaTeX}
\skillrow{Frameworks and Tools}{PyTorch, Docker, Z3, Semgrep, SonarQube, NuSMV/nuXmv, UPPAAL}
\skillrow{Methods}{Fuzzing, runtime verification, model checking, static analysis, and machine learning}

\end{document}
